\section{Real algebraic closure and canonical safe coordinates}
\label{sec:acf-geometry}

Throughout the following foundational geometric sections, the standing
symmetric binary AP and semantic homomorphic CCO hypotheses are in force.  We work in the
\emph{original} complete symmetric binary language.  In particular, all
colours, equivalences and quotients below are actual orbitals of the
homogeneous structure, not independently chosen labels.

For a finite real set \(A\subseteq M\), write
\(\operatorname{acl}_M(A)\) for its real model-theoretic algebraic
closure.  This set is finite by oligomorphy.  If \(S\subseteq
\mathcal D(M)\), say that \(S\) has \emph{strong closed amalgamation}
when every span of finite, induced, \(\operatorname{acl}_M\)-closed
structures has a strong amalgam in \(\Age(M)\) whose cross edges
between distinct vertices outside the common base have colours
in \(S\).  We neither assume strong amalgamation of arbitrary,
nonclosed finite structures nor require the amalgam itself to be
closed (it can be closed subsequently).

\begin{definition}[Rooted automorphism-transport ACF]
\label{def:rooted-acf}
For nonempty finite \(A_0\subseteq M\), a finite induced
\(W\subseteq M\) has \emph{rooted \(S\)-ACF over \(A_0\)} if it
contains \(\operatorname{acl}_M(A_0)\) and for every nonempty
\(Z\subseteq W\) such that
\(\col_W(x,y)\notin S\) whenever \(x,y\in Z\) are distinct,
and every selected \(z\in Z\), there exists
\(\alpha\in\Aut(W)\) satisfying
\[
 \alpha(Z)\subseteq\operatorname{acl}_M(A_0),
 \qquad \alpha(z)\in A_0.
\]
We also call such a \(Z\) an \emph{unsafe clique}.  The definition
requires one common \(W\) for all unsafe cliques, selected roots
and partial automorphisms of \(A_0\); it does not assert coherent
extensions.  In the AP-not-SAP situation the closure target is
essential: \(\operatorname{acl}_M(A_0)\) cannot be replaced in
general by \(A_0\).
\end{definition}

We first work in the \emph{finite-exact-normalised} case: no
nonidentity exact invariant equivalence has finite real classes.
Finite-exact normalisation will be used once in the common global
induction; no global ordinary EPPA result is assumed here.  Let \(F=\Sigma(\Id)\).  It can be a nonidentity, nonexact
thin packet equivalence, even in this normalised case.

\begin{lemma}[Finite exact intersections vanish]
\label{lem:acf-finite-exact-intersection}
If \(E\) is an exact invariant equivalence and \(a\) an actual
nonempty symmetric atomic orbital, then
\[
 E\cap a^2\subseteq F\quad\Longrightarrow\quad
 E\cap a^2=\Id.
\]
\end{lemma}
\begin{proof}
The intersection \(D=E\cap a^2\) is an exact symmetric
reflexive relation, by Lemma~\ref{lem:exact-product};
reflexivity uses transitivity of \(M\) and nonemptiness and
symmetry of \(a\).  If \(D\subseteq F\), its increasing positive
powers stabilise, because there are only finitely many original
orbitals.  The stable power is an exact invariant equivalence
contained in \(F\); it has finite classes.  Normalisation makes
it \(\Id\).  Since \(D\) is contained in its stable power,
\(D=\Id\).
\end{proof}

Let \(E_1,\ldots,E_r\) enumerate \emph{all} inclusion-minimal
infinite invariant equivalences.  By
Proposition~\ref{prop:minimalcut}, they are exact.  Define the
A32 correspondence relations
\[
 R_i=R_{E_i}
  =\bigcup\{a:\ a\text{ an atomic support (including equality), }
                         E_i\cap a^2=\Id\},
 \qquad C_i=\kappa(R_i).
 \tag{\(\star\)}\label{eq:acf-correspondence}
\]
Lemma~\ref{lem:global-correspondence} identifies every \(R_i\)
as an invariant equivalence and gives
\(F\le R_i\) and \(E_i\cap R_i=E_i\cap F\).
Every \(C_i>\Id\) is an exact equivalence by
Lemma~\ref{lem:packetinterval}, and
\([R_i:C_i]<\infty\).

\begin{lemma}[Concentration of unsafe pairs]
\label{lem:acf-unsafe-concentration}
Suppose \(S\) has strong closed amalgamation.  For any two
distinct minimal infinite cuts \(E_i,E_j\),
\[
   \mathcal D(M)\setminus S\ \subseteq\ R_i\cup R_j
   \quad\text{as unions of complete original orbitals}.
\]
Every finite unsafe clique is contained in one \(R_i\)-class
or in one \(R_j\)-class.
\end{lemma}
\begin{proof}
We first check an algebraicity test.  If \(xE_i u\) but
\(x\not F u\), then for every real \(v\),
\[
  x\in\operatorname{acl}_M(u,v)
             \quad\Longleftrightarrow\quad xR_i v.
  \tag{\(\dagger\)}\label{eq:acf-anchor-test}
\]
For the forward direction, the colour \(b=\col(x,u)\)
generates an infinite invariant equivalence inside the
minimal infinite \(E_i\); hence \(\Gamma(b)=E_i\).
The canonical real-hull theorem
(Theorem~\ref{thm:rh-hull}) implies that
\(E_i\cap\Gamma(\col(x,v))\) has finite classes, and thus is
contained in \(F\).  In particular
\(E_i\cap\col(x,v)^2\subseteq F\), which is \(\Id\) by
Lemma~\ref{lem:acf-finite-exact-intersection}.  This gives
\(xR_i v\).  The equality case \(x=v\) is immediate.
Conversely if \(xR_i v\), every element of the
\((u,v)\)-definable set
\[
       E_i[u]\cap\col(x,v)[v]
\]
is related to \(x\) by \(E_i\cap\col(x,v)^2=\Id\).
The set therefore consists of the single point \(x\);
in particular \(x\in\operatorname{acl}_M(u,v)\).
This proves~\eqref{eq:acf-anchor-test}.

Since \(E_i\cap E_j=\Id\) after normalisation (the meet
is an exact finite-class equivalence), every atom below
\(E_j\) has square below \(E_j\) and consequently belongs
to \(R_i\).  Thus \(E_j\le R_i\), and symmetrically
\(E_i\le R_j\).

Suppose an actual unsafe pair \(x,y\) lies in neither
\(R_i\) nor \(R_j\).  Choose \(u\in E_i[x]\setminus F[x]\)
and \(v\in E_j[y]\setminus F[y]\).  Since \(vE_jy\),
\(E_j\le R_i\), and \(x\not R_i y\), also \(x\not R_i v\).
The test gives \(x\notin\operatorname{acl}_M(u,v)\).
Likewise \(y\notin\operatorname{acl}_M(u,v)\).
Strongly \(S\)-amalgamate the finite closed structures
\(\operatorname{acl}_M(u,v,x)\) and
\(\operatorname{acl}_M(u,v,y)\) over
\(\operatorname{acl}_M(u,v)\).
Write \(x',y'\) for their distinct
off-base images and
\(a=\col(x,y)\) for the ORIGINAL
fine atomic colour, before amalgamation.
The amalgam preserves the COMPLETE
old pair colours from each
selected point to BOTH named
base points \(u,v\), not only the
displayed \(E_i\)- and \(E_j\)-incidences.

First \(x'E_i u\), while the old
pair \((u,y')\) has the same
colour as \((u,y)\).
Since \(uE_i x\), that latter
colour belongs to \(E_i;\{a\}\).
Transitivity \(E_i;E_i=E_i\)
therefore gives
\[
        \col(x',y')\in E_i;\{a\}.
\]
Independently the preserved
old colour of \(x'v\) equals
that of \(xv\), and \(vE_j y\),
so \(\col(x',v)\in\{a\};E_j\).
As \(vE_j y'\) and \(E_j\)
is transitive, we obtain
\[
        \col(x',y')\in\{a\};E_j.
\]
Same-sort invariant relations
commute in the ORIGINAL
symmetric binary algebra, so
the two restrictions give
\[
 \begin{aligned}
 \col(x',y')&
 \in (E_i;\{a\})\cap(E_j;\{a\})\\
 &=(E_i\cap E_j);\{a\}
   =\Id;\{a\}=\{a\},
 \end{aligned}
\]
where the intersection equality
is the FULL RAW exact meet-product
distributivity
(Lemma~\ref{lem:exact-product}),
and \(E_i\cap E_j=\Id\)
by finite-exact normalisation.
Thus the new cross edge has the
LITERAL original fine colour \(a\).
Since it is a NEW off-base
cross edge in an \(S\)-strong
amalgam, this colour must lie in
\(S\), contrary to the chosen
unsafe pair.  Hence every unsafe pair lies in
\(R_i\) or \(R_j\), and since these are unions of whole named
orbitals the same is true atomwise.

Finally, if every pair of a finite set is related by at least
one of two equivalences \(R_i,R_j\), the elementary
two-equivalence rook argument shows that the set lies
entirely inside one equivalence class of one of them:
a pair which is \(R_i\)-related but not \(R_j\)-related
forces every other vertex into its \(R_i\)-class.
\end{proof}

\begin{lemma}[Canonical exact MI coordinates]
\label{lem:acf-canonical-coordinates}
If \(r\ge2\), all \(C_i\) are infinite nonidentity exact
equivalences, and
\[
        \bigcap_iR_i=F,\qquad
        \bigcap_iC_i=\Id.
        \tag{\(\ddagger\)}\label{eq:acf-core-meets}
\]
For each \(i\), the actual exact quotient \(Q_i=M/C_i\)
has a meet-irreducible largest finite-class equivalence
\(F_i=R_i/C_i\), so
\[
     \operatorname{acl}_{Q_i}(U)=F_i[U]
               \quad\text{for every finite }U\subseteq Q_i.
               \tag{\(\heartsuit\)}\label{eq:acf-MI-closure}
\]
If \(S\) has strong closed amalgamation, every unsafe clique
\(Z\) satisfies, for some index \(i_0\),
\[
   \pi_j(Z)\subseteq F_j[\pi_j(z)]
              \quad(j\ne i_0,\ z\in Z).
        \tag{\(\clubsuit\)}\label{eq:acf-one-free}
\]
\end{lemma}
\begin{proof}
The intersection \(H=\bigcap_iR_i\) contains \(F\).
If it were strictly larger, it would have infinite classes
by maximality of \(F\), and hence contain some minimal
infinite \(E_j\).  But \(H\le R_j\) and
\(E_j\cap R_j=E_j\cap F\) is finite-class.  Thus \(H=F\).
The meet of the exact \(C_i\) is exact and contained in
\(F\); by normalisation it is \(\Id\).

For \(j\ne i\), \(E_j\le R_i\).  Since \(E_j\) is generated
by looped atoms and is exact,
\(E_j=\kappa(E_j)\le\kappa(R_i)=C_i\).  Thus \(C_i\)
is infinite and nonidentity, and every \(Q_i\) is a
strictly smaller original-language exact quotient.

Put \(R=R_i\) and \(E=E_i\).  If an invariant equivalence
\(D>R\), choose an atomic orbital \(a\subseteq D\setminus R\).
Then \(E\cap a^2\ne\Id\).  If \(E\cap D<E\), minimality of
the infinite \(E\) makes \(E\cap D\) finite-class, so
\(E\cap a^2\subseteq F\), contradicting
Lemma~\ref{lem:acf-finite-exact-intersection}.  Therefore
\(E\le D\); the least strict coarsening of \(R\) is \(E\vee R\).
Its relative index over \(R\) is infinite, since
\(E\cap R=E\cap F\) has finite classes.  Thus
\(\Sigma(R)=R\); Lemma~\ref{lem:packetinterval}
gives \(\Sigma(C_i)=R_i\).  Passing to the actual
exact quotient \(Q_i\), its greatest finite-class
equivalence is \(R_i/C_i\), and the upper interval is
meet-irreducible.  The unary-algebraicity assertion
\eqref{eq:acf-MI-closure} is the MI case of the
canonical real-hull criterion; equivalently it follows
from the absence of a pair of transverse infinite
directions in the quotient.

For every pair \(i\ne j\), Lemma~\ref{lem:acf-unsafe-concentration}
puts any unsafe clique wholly in an \(R_i\)-class or an
\(R_j\)-class.  Therefore at most one index \(i_0\)
can fail to contain the entire clique in a single
\(R_i\)-class.  For \(j\ne i_0\), passage through \(C_j\)
puts its projection into the finite \(R_j/C_j\) packet,
giving~\eqref{eq:acf-one-free}.
\end{proof}

For later reference we record the precise
closure pullback under FINITELY MANY
full exact quotient projections with
FINITE-CLASS common kernel. Let
\(K_1,\ldots,K_m\) be actual invariant
full exact equivalences and assume
\[
                     K=\bigcap_{j=1}^m K_j
\]
has FINITE real classes (in particular
this holds for \(K=\Id\)).
For every finite real \(A\),
\begin{equation}\label{eq:acf-closure-pullback}
 \operatorname{acl}_M(A)=
 \bigcap_{j=1}^m
 \pi_{K_j}^{-1}
       \bigl(\operatorname{acl}_{M/K_j}(\pi_{K_j}A)\bigr).
\end{equation}
For the left-to-right inclusion,
apply the exact quotient algebraic
projection lemma
\cref{lem:acf-acl-projects-exact}
to EACH original quotient
\(M\to M/K_j\): an algebraic
real point has algebraic projected
class. Conversely suppose the
\(m\) quotient projections of
\(x\) are all algebraic over
their corresponding finite
projected inputs. The
\(\Aut(M/A)\)-orbit of \(x\)
then projects into FINITELY
many quotient classes in each
of the \(m\) sources, hence
into finitely many joint
tuples of such classes.
Each joint tuple corresponds
to one \(K\)-class by actual
exact CRT, and that class has
FINITE size by hypothesis.
Thus the orbit of \(x\)
is finite, proving
\(x\in\operatorname{acl}_M(A)\).

The finite-class condition on
\(\bigcap_jK_j\) is ESSENTIAL.
For example in an infinite
source taking only
\(K_1=\Top\) makes the
right-hand side all of \(M\),
which cannot equal
\(\operatorname{acl}_M(A)\)
for finite \(A\).
We therefore never apply
\eqref{eq:acf-closure-pullback}
to an arbitrary family
of exact quotients without
checking the common kernel.

In particular, using \(K_i=C_i\),
\[
 \operatorname{acl}_M(A)=
   \left\{x:
     \pi_i(x)\in F_i[\pi_i(A)]
                    \text{ for all }i\right\}
 =\bigcap_iR_i[A].
 \tag{\(\spadesuit\)}\label{eq:acf-closure-bands}
\]

\begin{remark}[Avoiding arbitrary quotient safety]
The full-support safe-palette descent for the canonical
normalised factors will be proved after selecting maximal
exact cores in Lemma~\ref{lem:acf-quotient-safety} and
Corollary~\ref{cor:acf-hat-factors}.  It relies on a
\emph{harmless complementary exact split}.  The corresponding
statement for an arbitrary infinite exact quotient is false.
\end{remark}
